\begin{abstract}
\noindent
Lung cancer screening with low-dose computed tomography (CT) detects large
numbers of \emph{indeterminate} pulmonary nodules whose clinical fate is hard
to predict: some stay harmless for years, while others progress into aggressive
cancer. Today this risk stratification rests on the subjective visual judgement
of radiologists, which varies between readers and can lead either to
overtreatment of benign nodules or to dangerous delays for malignant ones. This
report asks whether an automated, \emph{intelligent} model can support that
decision by learning several biologically related properties of a nodule at
once.

\noindent
We propose \textit{MultiTaskCrossTalkNet}, a multi-task convolutional neural
network that jointly predicts a nodule's radiological subtype (ground-glass,
part-solid, solid) and its malignancy risk (benign vs.\ malignant). Instead of
rigidly sharing every layer, the network keeps one convolutional branch per task
and links them through learnable \emph{cross-talk} modules---a cross-stitch unit
and a co-attention unit---so each task can decide how much to borrow from the
other. We benchmark this design against single-task baselines that share no
information.

\noindent
All experiments use 508 nodule crops from the public LIDC-IDRI dataset,
converted to Hounsfield units, cropped around the annotated region of interest
and resized to $128\times128$ pixels. A third clinical task---two-year
progression---was part of the original design but had to be dropped, because the
progression labels in the clinical records could not be reliably linked to the
imaging annotations.

\noindent
Contrary to our hypothesis, cross-talk did \emph{not} consistently outperform
single-task learning under the conditions tested. The subtype task collapsed to
the majority ``solid'' class because of extreme class imbalance (97 of 101 test
samples were solid), and malignancy prediction gained nothing reliable from
information sharing. We attribute this to data quality and class distribution
rather than to the architecture itself, and we identify class-balanced sampling,
volumetric (3D) inputs, and a corrected progression linkage as the conditions
under which the multi-task hypothesis is most likely to pay off. We report this
honest negative result as a useful signal for others applying multi-task
learning to small, imbalanced medical-imaging datasets.

\end{abstract}